\documentclass{article}
\usepackage[utf8]{inputenc}
\usepackage[T1]{fontenc}
\usepackage{geometry}
\usepackage{graphicx}
\usepackage{float}
\usepackage{hyperref}
\usepackage[english,indonesian]{babel}

\title{Optimasi Pemrosesan Template LaTeX Menggunakan Pendekatan Berbasis AI}
\author{Academic Support Agent}
\date{\today}

\begin{document}
\maketitle

\begin{abstract}

\end{abstract}

\chapter{PENDAHULUAN}


\chapter{TINJAUAN PUSTAKA}


% Original content (commented out):
% \documentclass[11pt, a4paper, twocolumn]{article}
\usepackage[utf8]{inputenc}
\usepackage{multicol}
\usepackage{fancyhdr}
\usepackage{setspace}
\usepackage{indentfirst}
\usepackage{amsmath}
\usepackage{graphicx}
\usepackage{caption}
\captionsetup[table]{singlelinecheck=false}
\usepackage[subrefformat=parens,labelformat=parens]{subcaption}
\usepackage[a4paper]{geometry}
\geometry{top=1.95cm, bottom=1.8cm, left=2.4cm, right=2.4cm, headsep=0.5cm, headheight=1cm, 
            footskip=0in, footnotesep=0in, marginparwidth = 0pt,
            hoffset=0in, voffset=0cm}
\setlength{\parskip}{0cm}
\renewcommand{\baselinestretch}{1} 
\usepackage{hyperref}
\usepackage[backend=bibtex,style=numeric-comp,sorting=none,firstinits=true,maxbibnames=99]{biblatex}
\DeclareNameAlias{author}{last-first}
\bibliography{reference}


\pagestyle{fancy}
\renewcommand{\headrulewidth}{0pt}
\fancyhf{}
\chead{\footnotesize{ICLS2022 16th International Congress on Logistics and SCM Systems August 29-30, 2022 - Khon Kaen, Thailand.}}
\rfoot{\thepage}

\usepackage{sectsty} 
\sectionfont{\fontsize{12}{15}\selectfont}
\subsectionfont{\fontsize{12}{15}\selectfont}


\usepackage{lipsum}  

\makeatletter
\renewcommand{\maketitle}{\bgroup\setlength{\parindent}{0pt}
\begin{flushleft}
  \onehalfspacing
  \fontsize{20}{23}\selectfont
  \textbf{\@title} \\
  \hfill \break
  \fontsize{12}{15}\selectfont
  \@author
\end{flushleft}\egroup
}
\makeatother

\title{\LaTeX{} Template for 16th ICLS2022}
\author{%
        \textbf{Author One$^{1}$, Author Two$^{2,*}$}\\
        \fontsize{11}{14}\selectfont
        $^{1}$ First Author’s affiliation and address \\
        $^{2}$ Second Author’s affiliation and address \\
        Email: $^{1}$mail@domain.com, $^{2}$mail2@domain.com (Corresponding author)
        }

\begin{document}

\twocolumn[
  \begin{@twocolumnfalse}
    \maketitle

\noindent \textbf{Abstract.} An abstract is a shortened version of the paper and should contain all information necessary for the reader to determine: (1) what the objectives of the study were; (2) how the study was done; (3) what results were obtained; (4) the significance of the results. The typical length of an abstract is 150 – 300 words. \\

\noindent \textbf{Keywords:}  Keyword1, keyword2, keyword3, etc.

  \end{@twocolumnfalse}
  \vspace{1.5em}
]

% \begin{multicols}{2}

\section{Introduction}
\label{introduction}
We request authors to follow this guideline and format their manuscripts exactly the same as this document.

\section{Page Dimensions}
\label{pagedimension}
All material on all pages should fit within an area of A4 (21 x 29.7 cm), 1.95 cm from the top of the page and ending with 1.8 cm from the bottom. The left and right margins should both be 1.6 cm. This article demonstrates an example for you to prepare your article for submission to ICLS2022. Please strictly follow the guideline submit a compiled pdf version in the manuscript submission. You will be asked to submit the \LaTeX{} source files via email later. The maximum length is 6 pages including references.

\section{Main Text}
\label{maintext}
This section provides details for typesetting your manuscript according to the formatting guidelines set for ICLS2022 16th International Congress on Logistics and SCM Systems conference. Use regular font for typesetting of the main text in the document. The main text starts at the top of the page and continues in a two-column format. 

\subsection{Mathematics}
\label{mathematics}
Symbols and shorter expression can be placed within the text, e.g. $\epsilon \xrightarrow[]{}0$.  More complex expression should be placed in a center of a new line:

\begin{equation} \label{eq:01}
    E = mc^2
\end{equation}

Equations should be referred just like figures and tables, e.g. Eq. (\ref{eq:01}). Unless it is absolutely necessary, equation numbers should not have part to them, e.g. instead of using Eq. \ref{eq:01}(a) and Eq. \ref{eq:01}(b). Number them as Eq. (\ref{eq:01}) and Eq. (2).

\subsection{Figures/Table/Captions}
\label{caption}
Place tables/figures/images on the top of each column (see Table 1 and Fig. 1). Each figure or table must include a caption set in 11-point Times New Roman regular font, placed below the figure or table. The caption is left justified.

\begin{figure}[t]
    \centering
    \begin{subfigure}{0.4\linewidth}
    \includegraphics[width=\linewidth]{ICLSlogo_stroke.png}
    \caption{caption a}\label{fig:01a}
    \end{subfigure}
    \quad
    \begin{subfigure}{0.4\linewidth}
    \includegraphics[width=\linewidth]{ICLSlogo_stroke.png}
    \caption{caption b}\label{fig:01b}
    \end{subfigure}
    \caption{The figure is center-aligned and the caption of the figure is left-aligned.}
    \label{fig:01}
\end{figure}

Reference to the figure should follow the format “Fig. \ref{fig:01}”. Use “Figure \ref{fig:01}” instead if it is at the beginning of a sentence. Figure numbering and referencing should be done sequentially, e.g. Fig. \ref{fig:01}, Table \ref{tab:01}, etc. Use Fig. \ref{fig:01a}, Fig. \ref{fig:01b}, etc. for figures with multiple parts.

\begin{table}[t]
\caption{Table example}
\begin{tabular}{c c c c}
\hline
Heading & Column 1 & Column 2 & Column 3\\
\hline
Item 1 & 1 & 2 & 3 \\
Item 2 & 1 & 2 & 3 \\
\hline
\end{tabular}
\label{tab:01}
\end{table}

\section*{Acknowledgement}
\label{acknowledgement}
Funding supports should be acknowledged in this section.\\

\par

\section{In-text citation}
References should be arranged by the order in which they appear in the text. The example \LaTeX{} command for different in-text citation (please see the source code):
\begin{enumerate}
    \item Worasan et al. \cite{worasan2020hybrid}
    \item \citeauthor{worasan2020hybrid}
    \item \citeyear{worasan2020hybrid}
    \item Kusoncum et al. \cite{kusoncum2021heuristics}

\end{enumerate}

% \bibliographystyle{unsrt}
% \bibliography{reference}
\printbibliography



\end{document}


\end{document}